
In our research were able to highlight several directions that one could take.
We mainly highlight to directinos: \textbf{HF-PPAD} in  \textbf{PPAD}
and topological representations of directed graphs

\subsection{\textbf{HF-PPAD} in  \textbf{PPAD}}

In our sections \ref{sec:pc-plus-eol} and \ref{sec:sc-excess}, made references to
how Kleene logic can be incroportated with \textbf{PPAD} problems. For
the \textit{nD-StrongSperner} problems we observed that \textit{Kleene logic}
provides a natural extension and is equivalent to its non hazard counterpart.
For other problems such as \textit{StrongSperner}, we observed that this
is not clear by current complexity assumptions and indeed proving \textbf{PPAD}-hardness
would be essential for establishing the counting hardness of \textit{PureCircuit}.
But one could ask further questions such as, could we define
complexity classes such as \textbf{HF-PPAD} and could we show that
$\textbf{HF-PPAD} = \textbf{PPAD}$.
Along similar line, we can also ask the questions as to existence of problems
\textit{HF-nD-Counting-SS} and if they are parsimonious to their non hazard-free variants.
We know that the non counting version can be resolved parsimoniously
but we argue that it is not so simple when we omit panchromatic cubes as by doubling
the cubes, we may introduce more solutions.


Even more naturally, one can also ask further questions such as, how does hazard-freeness
affect \textbf{TFNP} as a whole. We conjecture that one can make
similar observations to the \textit{StrongTucker} problem \cite{aisenberg_2DTuckerPPA_2020},
in order to create similar classes such as \textbf{HF-PPA} as
it also uses a very similar colouring argument as the \textit{StrongSperner} argument.
Morevore we were able to observe how we can use \textit{Kleene logic} with
\textit{PLS} subproblems like \textit{Tarski}. We believe that this shift
in perspective may also give insights to other areas such as the circuit complexity \textit{Hazard-Free Circuits}.


\subsection{Parsimonious Topological Representations of \textbf{PPAD} problems }


From the previous sections we observed how to reduce \textit{SourceOrExcess}$(2,1)$ to a topological problem.
Just like in the \textit{EndOfLine} problem, we found it essential to reduce it to a \textit{StrongSperner} problem
with polynomial sides and very high dimensionality. We believe that one could abuse the \textit{Snake Embedding} methodology
that we proposed earlier but there are two main questions that need to be resolved:
\begin{enumerate}
    \item Will the number of solutions between foldings stay the same?
    \item Is there a natural way to define \textit{Counting-SS}?
\end{enumerate}
More interestingly, we made references of $\alpha-\textit{SUnbalanced}(\Delta_i, \Delta_o)$
and how they can affect the counting nature of each other. Like we mentioned previously, for
\textit{2D-StrongSperner} and \textit{1-MC}, we cannot guarantee parsimonious reductions, but
when we mapped our instances to a $3D$ space that was indeed possible. We therefore make the following conjecture:

\begin{conjecturebox}{\textit{nD-CountingSS} and \textit{Unbalanced}}{}
    There exist some polynomial computable function $f :\mathbb{N}^3 \to \mathbb{N}$, such that
    given $\alpha,\Delta_i, \Delta_o \in \mathbb{N}$ and $m= f(\alpha, \Delta_i, \Delta_o)$
    $$
        \textbf{PPAD}(\alpha\textit{-SUnbalanced}(\Delta_i, \Delta_o)) \subseteq
        \textbf{PPAD}(m\textit{D-CountingSS})
    $$
\end{conjecturebox}
What the above conjecture essentially describes is, given a number of known sources and the
maximum number of input/output vertices of graph, can we create some \textit{nD-CountingSS}
that is parsimonious to its graphical representation? Since we know that the problem \textit{Unbalanced}
can parsimoniously count every variant, can we also reduce \textit{Unbalanced} to a topological
problem? A reduction like that would be an important milestone for the usage topological \textit{PureCircuit} constructions
for enumerating solutions.